\documentclass[10pt,a4paper]{amsart}
\usepackage[margin=19mm]{geometry}
\usepackage{amsmath,amssymb,mathtools}
\usepackage[scr=rsfs,cal=euler]{mathalfa}
\usepackage{xfrac}
\usepackage[unicode,hidelinks,bookmarks=false]{hyperref}
\newcommand{\ie}{i.e.\ }
\newcommand{\cf}{cf.\ }
\newcommand{\resp}{respectively\ }
\newcommand{\Subsection}{\subsection}
\providecommand{\nobreakdash}{\nobreak\hbox{-}}
\setlength{\emergencystretch}{2em}
\multlinegap0pt
\usepackage{bm}
\usepackage{bbm}
\usepackage{dsfont}
\usepackage{xspace}
\usepackage{cancel}
\usepackage{mathtools}
\usepackage{amsthm}
\makeatletter
\@ifundefined{lemma}{\newtheorem{lemma}{Lemma}}{}
\makeatother
\makeatletter
\newcommand{\bB}{\mathbf{B}}
\newcommand{\bC}{\mathbf{C}}
\newcommand{\bbE}{\mathbb{E}}
\newcommand{\bh}{\mathbf{h}}
\newcommand{\calS}{\mathcal{S}}
\expandafter\ifx\csname vect\endcsname\relax
\newenvironment{vect}{\left[\begin{array}{c}}{\end{array}\right]}
\fi
\makeatother
\title[A Precise Performance Analysis of the Randomized Singular Va]{A Precise Performance Analysis of the Randomized Singular Value Decomposition}
\author{Danil Akhtiamov, Reza Ghane, Babak Hassibi}
\date{Source: arXiv:2510.06490v1 (CC BY 4.0)}
\begin{document}
\maketitle

\noindent\emph{Source and licence.} A Precise Performance Analysis of the Randomized Singular Value Decomposition. Version arXiv:2510.06490v1 (CC BY 4.0), distributed under the Creative Commons Attribution 4.0 International licence, as recorded in the arXiv licence metadata for this exact version and shown on the abstract page https://arxiv.org/abs/2510.06490v1. Full rights notice: licenses/arxiv-2510.06490-CC-BY-4.0.txt.\par\smallskip

\noindent\textit{Editorial scope.} Counted unit: the Lemma whose source label is \texttt{lm: HW\_bound} in the article (source0.tex lines 1044--1046, source label \texttt{lm: HW\_bound}), delivered together with the article's own complete original proof of it (source0.tex lines 1048--1157, one \texttt{proof} environment of 110 lines). No mathematics is written, repaired or re-derived: statement and proof are the article's own text line for line. Required context retained uncounted: eq: PO\_AO (lines 707--710). Every definition, display and label of those lines is retained unchanged. Omitted from this package and not redistributed: 1--706, 711--1043, 1047--1047, 1158--1773. Nothing inside a retained excerpt is omitted or altered: every retained body is delivered exactly as the article's lines are written, so no omission receipt of any kind accompanies this package. Citations: the retained excerpts carry no citation command at all, so no bibliography entry of the article is transcribed and no reference is redistributed. Reference adaptation: none was needed and none was made; every reference inside the delivered unit resolves inside it, so no unresolved reference or citation is produced. Printed slips kept literal: no printed word, symbol, hypothesis or number of the counted statement or of its proof was altered. Layout and class: the article's own class is \texttt{IEEEtran} with packages including amsmath, amssymb, amsfonts, amsthm, bm, algorithm2e, graphicx, textcomp, xcolor, subfig, xr, tikz, pgfplots, natbib; the wrapper is the lane's pinned \texttt{amsart} editorial wrapper at 10pt on A4, which restates the theorem-like environments the retained text needs and adds the article's own macro definitions that the retained text needs (\texttt{\textbackslash bB}, \texttt{\textbackslash bC}, \texttt{\textbackslash bbE}, \texttt{\textbackslash bh}, \texttt{\textbackslash calS}), with extra wrapper packages (bm, bbm, dsfont, xspace, cancel, mathtools). The delivered numbering is the unit's own. Assets: no retained span contains a figure environment, a graphics-inclusion command, a PDF-page inclusion, a table or a TikZ picture; no image file is copied into this package, the package declares no asset dependency and no image file is redistributed with it. Licence: version arXiv:2510.06490v1 of this article is distributed under the Creative Commons Attribution 4.0 International licence, as recorded in the arXiv licence metadata for this exact version and shown on the abstract page https://arxiv.org/abs/2510.06490v1. A full rights notice is delivered at licenses/arxiv-2510.06490-CC-BY-4.0.txt. Characters: every retained excerpt is pure ASCII, so the delivered engine, XeLaTeX with its default Latin Modern text font, reports no missing character.
\begin{align}
    \phi_i(\mathbf{g}, \mathbf{h}) :=  \max_{\mathbf{C}_i \in \calS_{\bC}} \min_{\mathbf{B}_i \in \calS_{\bB}} & \mathbf{C}_i^T \boldsymbol{\Sigma} \mathbf{g} \|\mathbf{B}_i\|   + \mathbf{B}_i^T \mathbf{h}\|\boldsymbol{\Sigma}\mathbf{C}_i\| \nonumber \\
    & + \sigma_i\mathbf{C}_{ii}- \frac{\|\mathbf{C}_i\|^2}{4} \label{eq: PO_AO}
\end{align}

\begin{lemma}\label{lm: HW_bound}
$$\mathbb{E} \phi_i(\mathbf{g}, \mathbf{h})   \sim \bbE \phi^{HW}_i(\mathbf{h}) $$
\end{lemma}

\begin{proof}
We will start with analyzing the objective without the expectation first. Continuing from the formulation provided in \eqref{eq: PO_AO}. By convexity-concavity of the objective in $\Phi_i$, we can swap the min and max and write
\begin{align*}
    \max_{\mathbf{C}_i \in \calS_{\bC}} \min_{\mathbf{B}_i \in \calS_{\bB}}  \mathbf{C}_i^T \boldsymbol{\Sigma} \mathbf{g} \|\mathbf{B}_i\| + \mathbf{B}_i^T \mathbf{h} \|\boldsymbol{\Sigma}\mathbf{C}_i\| + \sigma_i\mathbf{C}_{ii}- \frac{\|\mathbf{C}_i\|^2}{4}
\end{align*}
We let $\bB_i = \beta \hat{\bB}_i$ where $\|\hat{\bB}_i\|_2 = 1$. We perform the optimization over the direction, $\hat{\bB_i}$, which yields the following equivalent optimization:
\begin{align*}
     \max_{\mathbf{C}_i \in \calS_{\bC}} \min_{\beta \ge 0} \mathbf{C}_i^T \boldsymbol{\Sigma} \mathbf{g} \beta - \beta \|\mathbf{h}\| \|\boldsymbol{\Sigma}\mathbf{C}_i\| + \sigma_i\mathbf{C}_{ii}- \frac{\|\mathbf{C}_i\|^2}{4} 
\end{align*}
Now we use the square-root trick $\sqrt{x}:= \min_{\tau \ge 0} \frac{1}{\tau} + \tau x^2$ to obtain the following optimization
\begin{align*}
     \max_{\mathbf{C}_i \in \calS_{\bC}} \min_{\beta \ge 0} \max_{\tau \ge 0} & \;\mathbf{C}_i^T \boldsymbol{\Sigma} \mathbf{g} \beta - \frac{\beta  \|\mathbf{h}\|}{2} \left(\tau\|\boldsymbol{\Sigma}\mathbf{C}_i\|^2 + \frac{1}{\tau}\right) \\ &+ \sigma_i\mathbf{C}_{ii}- \frac{\|\mathbf{C}_i\|^2}{4}
\end{align*}
Rearranging terms yields
\begin{align}
    \max_{\mathbf{C}_i \in \calS_{\bC}}  \min_{\beta \ge 0} \max_{\tau \ge 0} - \frac{\beta\|\mathbf{h}\|}{2\tau} + \mathbf{C}_i^T \left(\beta \boldsymbol{\Sigma} \mathbf{g} + \sigma_i\mathbf{e}_i\right) \nonumber \\
    \quad - \frac{\beta\tau  \|\mathbf{h}\|}{2} \|\boldsymbol{\Sigma}\mathbf{C}_i\|^2 - \frac{\|\mathbf{C}_i\|^2}{4} \label{opt: scal_not}
\end{align}

Now we observe that the objective in \eqref{opt: scal_not} is convex-concave with respect to $\beta$ and $(\tau, \bC_i)$ respectively. Hence using Sion's minimax theorem, we may swap the order of the minimization and maximizations to obtain
\begin{align*}
    \min_{\beta \ge 0} \max_{\tau \ge 0} \max_{\mathbf{C}_i \in \calS_{\bC}}  - \frac{\beta\|\mathbf{h}\|}{2\tau} + \mathbf{C}_i^T \left(\beta \boldsymbol{\Sigma} \mathbf{g} + \sigma_i\mathbf{e}_i\right) \\
    \quad - \frac{\beta\tau  \|\mathbf{h}\|}{2} \|\boldsymbol{\Sigma}\mathbf{C}_i\|^2 - \frac{\|\mathbf{C}_i\|^2}{4}
\end{align*}
Performing the quadratic optimization over $\bC_i$ yields
\begin{align}\label{eq: obj_before_HW}
& \min_{\beta \ge 0} \max_{\tau \ge 0}  - \frac{\beta\|\mathbf{h}\|}{2\tau} \nonumber \\ 
& + \Bigl(\beta \boldsymbol{\Sigma} \mathbf{g} + \sigma_i\mathbf{e}_i\Bigr)^T\Bigl(\mathbf{I} + 2 \beta\tau \|\mathbf{h}\|\boldsymbol{\Sigma}^2\Bigr)^{-1} \Bigl(\beta \boldsymbol{\Sigma} \mathbf{g} + \sigma_i\mathbf{e}_i\Bigr) 
\end{align}

% Let $\mathbf{M} = \Bigl(\mathbf{I} + 2\beta\tau \|\mathbf{h}\|\boldsymbol{\Sigma}^2\Bigr)^{-1}$. Expanding the quadratic form:
% \begin{align*}
% &\Bigl(\beta \boldsymbol{\Sigma} \mathbf{g} + \sigma_i\mathbf{e}_i\Bigr)^T\mathbf{M}\Bigl(\beta \boldsymbol{\Sigma} \mathbf{g} + \sigma_i\mathbf{e}_i\Bigr) = \\
% & \beta^2 \mathbf{g}^T \boldsymbol{\Sigma} \mathbf{M} \boldsymbol{\Sigma} \mathbf{g} + 2\beta\sigma_i \mathbf{g}^T \boldsymbol{\Sigma} \mathbf{M} \mathbf{e}_i + \sigma_i^2 \mathbf{e}_i^T \mathbf{M} \mathbf{e}_i
% \end{align*}

% Since $\mathbf{g} \sim \mathcal{N}(0, \mathbf{I})$, by Hanson-Wright inequality:
% \begin{align*}
% &\mathbb{P}\Bigl(\Bigl|\mathbf{g}^T \boldsymbol{\Sigma} \mathbf{M} \boldsymbol{\Sigma} \mathbf{g} - \mathrm{tr}(\boldsymbol{\Sigma} \mathbf{M} \boldsymbol{\Sigma})\Bigr| > t\Bigr) \\
% &\leq 2\exp\left(-\frac{1}{2}\min\left(\frac{t^2}{4\|\boldsymbol{\Sigma} \mathbf{M} \boldsymbol{\Sigma}\|_F^2}, \frac{t}{2\|\boldsymbol{\Sigma} \mathbf{M} \boldsymbol{\Sigma}\|_{\text{op}}}\right)\right) \nonumber
% \end{align*}

% For the cross term, since $\mathbf{g} \sim \mathcal{N}(0, \mathbf{I})$ and $\boldsymbol{\Sigma} \mathbf{M} \mathbf{e}_i$ is deterministic:
% \begin{align*}
% \mathbb{P}\Bigl(\Bigl|\mathbf{g}^T \boldsymbol{\Sigma} \mathbf{M} \mathbf{e}_i\Bigr| > t\Bigr) \leq 2\exp\left(-\frac{t^2}{2\|\boldsymbol{\Sigma} \mathbf{M} \mathbf{e}_i\|^2}\right)
% \end{align*}

As $m \to \infty$, due to the Hanson-Wright inequality, \eqref{eq: obj_before_HW} converges in probability to:
\begin{align*}
\min_{\beta \ge 0} \max_{\tau \ge 0} &- \frac{\beta\|\mathbf{h}\|}{2\tau} +\beta^2 \mathrm{tr} \left(\boldsymbol{\Sigma}^2\Bigl(\mathbf{I} + 2 \beta\tau \|\mathbf{h}\|\boldsymbol{\Sigma}^2\Bigr)^{-1}\right) \nonumber\\
&+ \frac{\sigma_i^2}{1 + 2\beta\tau\|\mathbf{h}\|\sigma_i^2} \nonumber
\end{align*}
Which can be rewritten as:
\begin{align}\label{eq: right_after_HW}
\min_{\beta \ge 0} \max_{\tau \ge 0} - \frac{\beta\|\mathbf{h}\|}{2\tau} +  \sum_j \frac{\sigma_j^2(\beta^2 + \delta_{ij})}{1 + 2\beta\tau\|\mathbf{h}\|\sigma_j^2} 
\end{align}


Setting derivatives by $\tau$ and $\beta$ to zero:
\begin{align*}
&\frac{\beta\|\mathbf{h}\|}{2\tau^2} -  \sum_j \frac{2\sigma_j^4\beta\|\mathbf{h}\|(\beta^2 + \delta_{ij})}{(1 + 2\beta\tau\|\mathbf{h}\|\sigma_j^2)^2} = 0, \\
&- \frac{\|\mathbf{h}\|}{2\tau} + \sum_j \sigma_j^2 \frac{2\beta(1 + 2\beta\tau\|\mathbf{h}\|\sigma_j^2)}{(1 + 2\beta\tau\|\mathbf{h}\|\sigma_j^2)^2} \nonumber\\
&\quad - \sum_j \sigma_j^2\frac{2\tau \|\mathbf{h}\|\sigma_j^2(\beta^2 + \delta_{ij})}{(1 + 2\beta\tau\|\mathbf{h}\|\sigma_j^2)^2} = 0 \nonumber
\end{align*}

Simplifying both equations we arrive at:
\begin{align*}
\frac{1}{2\tau^2} &=  \sum_j \frac{2\sigma_j^4(\beta^2 + \delta_{ij})}{(1 + 2\beta\tau\|\mathbf{h}\|\sigma_j^2)^2} \\
\frac{\|\mathbf{h}\|}{2\tau}  &= \sum_j \sigma_j^2 \frac{2\beta + 2\beta^2\tau\|\mathbf{h}\|\sigma_j^2 - 2\tau \|\mathbf{h}\|\sigma_j^2 \delta_{ij}}{(1 + 2\theta\|\mathbf{h}\|\sigma_j^2)^2} \nonumber
\end{align*}

Introducing $\tilde{\theta}= \beta\tau$ we get:
\begin{align*}
1 &=  \sum_j \frac{4\sigma_j^4(\tilde{\theta}^2 + \delta_{ij}\tau^2)}{(1 + 2\tilde{\theta}\|\mathbf{h}\|\sigma_j^2)^2} \\
\|\mathbf{h}\| &= \sum_j 4 \sigma_j^2 \frac{\tilde{\theta} + \tilde{\theta}^2\|\mathbf{h}\|\sigma_j^2 - \tau^2 \|\mathbf{h}\|\sigma_j^2 \delta_{ij}}{(1 + 2\tilde{\theta}\|\mathbf{h}\|\sigma_j^2)^2} \nonumber
\end{align*}

Replacing the penultimate equation above by the the sum of the last equation with the penultimate one multiplied by $\|\mathbf{h}\|$ and canceling out $\tilde{\theta}$ we arrive at:
% \begin{align*}
% 2\|\mathbf{h}\| &= \sum_j 4 \sigma_j^2 \frac{\tilde{\theta} + 2\tilde{\theta}^2\|\mathbf{h}\|\sigma_j^2}{(1 + 2\tilde{\theta}\|\mathbf{h}\|\sigma_j^2)^2}\nonumber
% \end{align*}
%  leads to the following equation:
\begin{align}\label{eq: theta}
\sum_j 2 \sigma_j^2 \frac{\tilde{\theta}}{1 + 2\tilde{\theta}\|\mathbf{h}\|\sigma_j^2} = \|\mathbf{h}\| 
\end{align}
The value of \eqref{eq: right_after_HW} simplifies using \eqref{eq: theta} as:
\begin{align*}
&- \frac{\tilde{\theta}\|\mathbf{h}\|}{2\tau^2} +  \sum_j \frac{\sigma_j^2(\beta^2 + \delta_{ij})}{1 + 2\tilde{\theta}\|\mathbf{h}\|\sigma_j^2} = \frac{\sigma_i^2}{1 + 2\tilde{\theta}\|\mathbf{h}\|\sigma_i^2}
\end{align*}
The value of \eqref{eq: right_after_HW} simplifies as follows using \eqref{eq: theta}:
\begin{align*}
&- \frac{\tilde{\theta}\|\mathbf{h}\|}{2\tau^2} +  \sum_j \frac{\sigma_j^2(\beta^2 + \delta_{ij})}{1 + 2\tilde{\theta}\|\mathbf{h}\|\sigma_j^2} = - \frac{\tilde{\theta}\|\mathbf{h}\|}{2\tau^2} +  \sum_j \frac{\sigma_j^2(\frac{\tilde{\theta}^2}{\tau^2} + \delta_{ij})}{1 + 2\tilde{\theta}\|\mathbf{h}\|\sigma_j^2 } \\
&=   \frac{\tilde{\theta}}{2\tau^2} \left( \sum_j \frac{2\tilde{\theta} \sigma_j^2}{1 + 2\tilde{\theta}\|\mathbf{h}\|\sigma_j^2 } -\|\mathbf{h}\| \right) +   \frac{\sigma_i^2}{1 + 2\tilde{\theta}\|\mathbf{h}\|\sigma_i^2}\\ & = \frac{\sigma_i^2}{1 + 2\tilde{\theta}\|\mathbf{h}\|\sigma_i^2}
\end{align*}
We would also like to slightly simplify \eqref{eq: theta}
\begin{align*}
& \|\mathbf{h}\|^2 = \sum_j 2 \sigma_j^2 \frac{\tilde{\theta}\|\mathbf{h}\| }{1 + 2\tilde{\theta}\|\mathbf{h}\|\sigma_j^2} = \\
&  \sum_j 1 - \frac{1}{1 + 2\tilde{\theta}\|\mathbf{h}\|\sigma_j^2} = m - \sum_{j=1}^m \frac{1}{1 + 2\tilde{\theta}\|\mathbf{h}\|\sigma_j^2}
\end{align*}
% Thus
% Also note that \eqref{eq: theta} can be simplified as: 
% \begin{align*}
% m - \|\mathbf{h}\|^2 = \sum_{j=1}^m \frac{1}{1 + 2\tilde{\theta}\|\mathbf{h}\|\sigma_j^2}
% \end{align*}
Denoting $\theta_{\bh} := 2\theta \|\mathbf{h}\|$ finally yields the desired expression $\frac{\sigma_i^2}{1 +\theta_{\bh} \sigma_i^2}$ for \eqref{eq: obj_before_HW}, where $\theta_{\bh}$ is defined through
\begin{align*}
\sum_{j=1}^m \frac{1}{1 + \theta_\bh \sigma_j^2} = m - \|\mathbf{h}\|^2 
\end{align*}

\end{proof}

\end{document}
